%=====================================================================
\chapter{Product Variations, Options and Uploads}\label{ch:variations}
%=====================================================================
\locator{2}{Part II --- Standing the Shop Up}

\begin{outcomes}
By the end of this chapter you will be able to:
\begin{tight}
  \item \textbf{Distinguish} an option from a variant, and \textbf{say} which
        of the two a given requirement needs.
  \item \textbf{Explain} why options do not store the combinations they imply,
        and \textbf{state} what that makes impossible.
  \item \textbf{Configure} option sets on a product and \textbf{predict} the
        number of rows each one creates.
  \item \textbf{Describe} what a shop must check before accepting an uploaded
        file, and \textbf{identify} which of those checks a client can defeat.
  \item \textbf{Test} your own shop's upload handling and \textbf{report} what
        it accepts.
\end{tight}
\end{outcomes}

\begin{prereq}
\chref{catalogue}: the eighteen product tables, and the habit of looking at
what a screen wrote. Your shop should still have its demonstration catalogue;
\texttt{seed/catalogue.py --restore} puts it back if not.
\end{prereq}

\begin{keyterms}
option $\bullet$ option value $\bullet$ variant $\bullet$ SKU $\bullet$
modifier $\bullet$ price prefix $\bullet$ combinatorial explosion $\bullet$
stock keeping $\bullet$ required option $\bullet$ file upload $\bullet$
allow-list $\bullet$ MIME type $\bullet$ content sniffing $\bullet$ stored
cross-site scripting
\end{keyterms}

\section{Two Ways to Say ``Comes in Three Colours''}

A shirt in three colours and four sizes. Twelve things a customer might buy.
How many rows is that?

The honest answer is that it depends on a decision nobody will prompt you to
make, and that the two answers have different consequences for the rest of the
shop's life.

\begin{definitionbox}[Option and Variant]
An \term{option} is a choice attached to one product, whose values
\emph{modify} it. The shirt is one product; ``Large'' adds Rs.~200 and
``Red'' adds nothing. The twelve combinations are \textbf{not stored
anywhere} --- they are computed when a customer picks.

\smallskip
A \term{variant} is a product in its own right, linked to a master. ``Red
Large Shirt'' is a row in \texttt{oc\_product}, with
\texttt{master\_id} pointing at the parent, its own model number, its own
stock and its own price. Twelve combinations are twelve rows.

\smallskip
OpenCart~4 supports both --- \texttt{oc\_product} carries
\texttt{master\_id}, \texttt{variant} and \texttt{override} columns alongside
the older option tables. The admin panel shows options first, so options are
what most shops end up with, including shops that needed variants.
\end{definitionbox}

\subsection{What options actually store}

Two tables do the work, and a third names the choices globally.

\rowcolors{2}{white}{lightbg}
\begin{longtable}{@{}L{5.4cm}L{9.6cm}@{}}
\toprule
\rowcolor{primary!15}
\textbf{Table} & \textbf{What a row means} \\
\midrule
\endhead
\texttt{oc\_option} & A choice that exists in the shop at all --- ``Size'',
``Colour'' --- with its \texttt{type}: \texttt{select}, \texttt{radio},
\texttt{checkbox}, \texttt{text}, \texttt{textarea}, \texttt{date},
\texttt{time}, \texttt{datetime}, \texttt{file} \\
\texttt{oc\_option\_value} & One possible value of it: ``Large''. Shared
across every product that uses the option \\
\texttt{oc\_product\_option} & \emph{This} product offers \emph{that} option,
and whether it is \texttt{required} \\
\texttt{oc\_product\_option\_value} & This product offers that value, and
what choosing it does: \texttt{price} with a \texttt{price\_prefix} of
\texttt{+} or \texttt{-}, plus \texttt{quantity}, \texttt{weight} and
\texttt{points} adjustments \\
\bottomrule
\end{longtable}

\begin{alertbox}[What options make impossible]
Look at where \texttt{quantity} lives: on
\texttt{oc\_product\_option\_value}, which is \emph{per option value}, not per
combination.

\smallskip
So the shop can know it has 40 Large shirts and 30 Red shirts. It
\textbf{cannot} know how many Red Large shirts it has, because there is no row
for that --- and no amount of configuration will add one. Three consequences
follow, and all three are discovered late:

\begin{tight}
  \item \textbf{No stock per combination.} You will oversell the popular
        combination and sit on the unpopular one.
  \item \textbf{No SKU or barcode per combination}, so the shop cannot speak
        to a warehouse system that works in SKUs --- which is all of them.
  \item \textbf{No price history or reporting per combination.} ``How many
        Red Large did we sell?'' has no answer.
\end{tight}

\smallskip
If any of those three matter, you need variants, and you need to know that
\emph{before} the catalogue is built. Migrating two thousand option-based
products to variants is not a setting; it is a project.
\end{alertbox}

\begin{conceptbox}[Which one does this requirement need?]
A test that settles most cases in one question: \textbf{does the shop need to
count, price or identify the combination on its own?}

\smallskip
\textbf{Options are right} when the choice does not change what you hold in
stock: an engraving message, a delivery date, a gift-wrap tick, a cable length
that you cut to order. Also when the combination genuinely shares stock ---
which is rarer than people assume.

\smallskip
\textbf{Variants are right} whenever the combination is a thing in a box on a
shelf. Clothing, shoes, anything with a barcode. If a warehouse person could
pick up exactly one combination and scan it, it is a variant.

\smallskip
\textbf{The common mistake} is to choose options because the admin panel
offers them first and they are quicker to set up, and to discover in month
nine --- during the first real stock-take --- that the shop has never known
what it holds.
\end{conceptbox}

\section{What Each Choice Costs}

Options are cheap to store and variants are not. That much is obvious. What is
worth measuring is \emph{how fast} the difference grows, because it decides
whether a catalogue is manageable.

\begin{measured}[Options Against Variants]
Produced by \texttt{code/m07\_variations.py} on the machine in Appendix~A.
Option sets are attached to one product and the product page is requested seven
times, median taken; the variant column is the number of
\texttt{oc\_product} rows the same catalogue would need.

\rowcolors{2}{white}{lightbg}
% Six columns inside a tcolorbox: 10.4 cm of content plus five 0.42 cm gaps.
% Seven columns overflowed by 12.3 pt, and "option rows" was the column the
% prose could carry instead.
\begin{longtable}{@{}L{1.9cm}L{1.7cm}L{1.7cm}L{1.6cm}L{1.4cm}L{2.1cm}@{}}
\toprule
\rowcolor{primary!15}
\textbf{Options} & \textbf{Combin.} & \textbf{Value rows} & \textbf{Page} &
\textbf{KB} & \textbf{As variants} \\
\midrule
\endhead
3 & 3 & 3 & 87\,ms & 30.0 & 4 rows \\
3$\times$4 & 12 & 7 & 104\,ms & 30.7 & 13 rows \\
3$\times$4$\times$5 & 60 & 12 & 101\,ms & 31.3 & 61 rows \\
6$\times$8$\times$5 & \textbf{240} & \textbf{19} & 103\,ms & 31.3 &
\textbf{241 rows} \\
\bottomrule
\end{longtable}

\smallskip
\textbf{Options grow by addition; variants grow by multiplication.} Three
option sets of 6, 8 and 5 values are \textbf{19 rows} --- $6+8+5$ --- and
express 240 combinations. The same catalogue as variants is \textbf{241
product rows}, each with its own description row, its own store row and its
own category rows. Add a sixth colour and options gain one row while variants
gain forty.

\smallskip
\textbf{And the product page does not care.} From 3 combinations to 240, the
page stayed at about 100\,ms and grew by \textbf{1.3 kilobytes} --- because the
page renders the \emph{choices}, not the combinations. Nineteen values in three
dropdowns, whatever the 240 products those values imply.

\smallskip
\textbf{So the cost of variants is not page speed; it is catalogue
management.} Two hundred and forty-one rows that must each be priced, stocked,
described, photographed and kept in step. That is the real bill, it is paid by
a person rather than a server, and it is the reason shops that should have
chosen variants did not.
\end{measured}

\section{Uploads, and What Your Shop Will Accept}

One of the nine option types is \texttt{file}. Switch it on and a customer can
send you a file --- a logo to print, a document to quote against, a photograph
to frame.

You have now accepted arbitrary bytes from a stranger, and everything that
follows is about not regretting it.

\begin{examplebox}[What OpenCart checks, in order]
From \texttt{catalog/controller/tool/upload.php} in your own container. Read
it; it is forty lines.

\begin{tightnum}
  \item \textbf{The filename is sanitised} --- everything outside
        \texttt{[a-zA-Z0-9.\-\textvisiblespace+]} is stripped, then
        \texttt{basename()} removes any path. This defeats
        \texttt{../../etc/passwd} and it is solid.
  \item \textbf{The length is checked}, 3 to 64 characters.
  \item \textbf{The extension is checked against an allow-list}
        (\texttt{config\_file\_ext\_allowed}). An allow-list is the right
        shape: anything not named is refused.
  \item \textbf{The MIME type is checked against a second allow-list}
        (\texttt{config\_file\_mime\_allowed}).
  \item \textbf{The stored name is randomised} --- the file is saved as the
        name plus a 32-character token, so it cannot be guessed, and the real
        name is kept in the database.
\end{tightnum}
\end{examplebox}

\begin{alertbox}[Check 4 is worth almost nothing, and it is not a bug in OpenCart]
The MIME type it compares is
\texttt{\$\_FILES['file']['type']} --- and that value is \textbf{sent by the
browser}. It is a claim made by the client about its own file. Anyone using
anything other than a browser sets it to whatever passes.

\smallskip
The default allow-list makes this worse by including
\texttt{application/octet-stream}, which means ``unidentified binary'' ---
so the honest answer for almost any file is already on the list.

\smallskip
\textbf{This is not a criticism of OpenCart so much as a lesson about where a
check must run.} A check performed on data the attacker controls is not a
check. The extension allow-list is real because the server derives it from the
name it has already sanitised. The MIME check is advisory, and the only
server-side equivalent is to inspect the file's actual contents ---
\texttt{finfo\_file()} in PHP, which reads the magic bytes.
\end{alertbox}

\begin{alertbox}[Two defaults on your shop, which you can verify in five minutes]
Both of these are true of the stock install in \chref{stack}, and both are
measured in \chref{security}.

\smallskip
\textbf{The allow-list includes \texttt{svg}.} An SVG is not a picture; it is
an XML document, and it may contain \texttt{<script>}. A browser asked to
display one from your origin will run it. The default list also carries
\texttt{eps} and \texttt{ps}, which are PostScript --- a programming language
--- and \texttt{msi}, \texttt{cab} and \texttt{rar}.

\smallskip
\textbf{Uploads are served over HTTP.} They land in
\texttt{system/storage/upload/}, which is \emph{inside} the web root. Requesting
a file there returns it: a text file comes back as \texttt{text/plain} and an
SVG as \texttt{image/svg+xml}. The directory cannot be listed --- there is an
empty \texttt{index.html} --- and names carry a 32-character token, so this is
not trivially exploitable. But the protection is secrecy of the name, not
inaccessibility of the directory, and Apache is configured with
\texttt{AllowOverride None}, so the \texttt{.htaccess} OpenCart ships cannot
change it.

\smallskip
\textbf{What to do about it} is \chref{security}'s subject. The short version:
remove \texttt{svg} from the list unless you need it, serve uploads through a
script that sets \texttt{Content-Disposition: attachment} rather than from the
filesystem, and keep the storage directory outside the web root.
\end{alertbox}

\begin{conceptbox}[Two size limits, and the smaller one wins]
\texttt{config\_file\_max\_size} is a setting in the admin panel. PHP has its
own, \texttt{upload\_max\_filesize}, in \texttt{code/shop/php.ini}, and this
course sets it to 8~MB.

\smallskip
If the admin setting is larger, it has no effect: PHP refuses the upload
before any OpenCart code runs, and the error the customer sees is unhelpful.
Whenever a limit appears not to work, look for the second limit --- there is
nearly always a second limit, in the web server, the runtime or the proxy, and
the smallest one is the one in force.
\end{conceptbox}

\begin{pitfall}
\begin{tight}
  \item \textbf{Choosing options when the requirement needs variants.} The
        symptom arrives at the first stock-take, and the fix is a migration.
  \item \textbf{Believing option \texttt{quantity} tracks combinations.} It
        tracks one value. Red Large and Red Small share Red's stock.
  \item \textbf{Trusting the uploaded MIME type.} The client supplies it.
  \item \textbf{Allowing \texttt{svg} without thinking about it.} It is script
        that renders as a picture.
  \item \textbf{Setting a size limit in one place only.} The smaller of the
        admin setting and \texttt{upload\_max\_filesize} wins.
  \item \textbf{Making every option \texttt{required}.} The customer cannot
        buy without choosing, and a required text option with no sensible
        default costs sales.
  \item \textbf{Forgetting that option values are global.} Editing ``Large''
        edits it for every product that uses it.
\end{tight}
\end{pitfall}

\begin{lab}[Options, variants, and what your shop will swallow]
\begin{lstlisting}[language=bash]
# --- 1. Look at a product that already has every option type ---------
# Product 42 in the demo catalogue carries all nine. Open it in the
# admin panel: Catalog > Products > edit > Option.
cd code
docker compose exec db mariadb -ushop -pshoppw opencart

# --- 2. Count what one option set costs ------------------------------
# Add a Size option with three values, through the admin panel. Then
# count the rows it created. Predict the number before you look.
#   SELECT COUNT(*) FROM oc_product_option       WHERE product_id = 42;
#   SELECT COUNT(*) FROM oc_product_option_value WHERE product_id = 42;

# --- 3. Find the row that does not exist -----------------------------
# Set stock on two option values. Then try to write a query that
# answers "how many Red Large do we hold?". You cannot. Write down,
# in one sentence, why -- that sentence is the whole of section 7.1.

# --- 4. Measure options against variants -----------------------------
python m07_variations.py

# --- 5. Read the upload controller -----------------------------------
# Forty lines. Find the five checks, and find the one the client
# controls.
docker compose exec shop \
  cat /var/www/html/catalog/controller/tool/upload.php

# --- 6. Find out what your shop allows -------------------------------
docker compose exec db mariadb -ushop -pshoppw opencart -e \
  "SELECT \`key\`,\`value\` FROM oc_setting \
    WHERE \`key\` LIKE 'config_file%';"

# --- 7. Prove the upload directory is served -------------------------
# Put a file there yourself and fetch it. This is YOUR container; see
# the scope statement in Appendix A before you try this anywhere else.
docker compose exec shop sh -c \
  'echo hello > /var/www/html/system/storage/upload/probe.txt'
curl -i http://localhost:8090/system/storage/upload/probe.txt
docker compose exec shop rm /var/www/html/system/storage/upload/probe.txt

# --- 8. Decide, and write it down ------------------------------------
# For the shop you described in Chapter 4: options or variants? One
# paragraph, with the stock question answered explicitly. Chapter 23
# marks you on whether the catalogue matches this decision.
\end{lstlisting}

\smallskip
\textbf{What to notice.} Step 3 is the one that teaches. Everybody believes
option stock tracks combinations until they try to write the query, and the
minute spent failing to write it is worth more than this chapter's prose.
\end{lab}

\begin{checkpoint}
\begin{tightnum}
  \item A shop sells mugs with a customer-supplied photograph printed on them,
        in two sizes. Which choices are options and which are variants, and
        why? Say what the shop must hold in stock.
  \item A product has option sets of 4, 3 and 2 values. How many rows in
        \texttt{oc\_product\_option\_value}, and how many
        \texttt{oc\_product} rows would the same catalogue need as variants?
  \item An attacker uploads a file named \texttt{logo.svg} containing
        \texttt{<script>}, sending \texttt{Content-Type:} \texttt{image/png}.
        Which of OpenCart's five checks does it pass, which does it fail, and
        what would have to be true for this to harm a visitor?
\end{tightnum}
\end{checkpoint}

\begin{chaptersummary}
\begin{tight}
  \item An option modifies one product; a variant is a product of its own with
        \texttt{master\_id} pointing at a master. OpenCart~4 supports both and
        shows options first.
  \item Options store choices, not combinations. Stock, price and weight sit on
        the option \emph{value}, so there is no row for ``Red Large'' and no
        query can produce one.
  \item That makes three things impossible: stock per combination, a SKU per
        combination, and reporting per combination. If any matters, use
        variants, and decide before the catalogue is built.
  \item The test: could a warehouse person pick up exactly this combination and
        scan it? Then it is a variant.
  \item Measured: option sets of 6, 8 and 5 values are 19 rows and express 240
        combinations; the same catalogue as variants is 241 product rows.
        Options grow by addition, variants by multiplication --- and the page
        time is unaffected either way, because the page renders choices.
  \item The real cost of variants is catalogue management by a person, not
        server time.
  \item A \texttt{file} option accepts arbitrary bytes from a stranger.
        OpenCart sanitises the name, checks the extension against an allow-list
        and randomises the stored name --- all sound. Its MIME check is not,
        because the client supplies the MIME type.
  \item On the stock install \texttt{svg} is allowed and the upload directory
        is served over HTTP. The protection is an unguessable name, not an
        unreachable directory. \chref{security} fixes both.
\end{tight}
\end{chaptersummary}

\begin{reviewq}
  \item \textbf{(MCQ)} Stock quantity for an option lives on:
        (a)~\texttt{oc\_product}; (b)~\texttt{oc\_product\_option};
        (c)~\texttt{oc\_product\_option\_value}; (d)~a combination table.
  \item \textbf{(MCQ)} Three option sets of 5, 4 and 3 values imply how many
        combinations, and cost how many \texttt{oc\_product\_option\_value}
        rows? (a)~12 and 12; (b)~60 and 12; (c)~60 and 60; (d)~12 and 60.
  \item \textbf{(Short)} Give the one-question test for options against
        variants, and apply it to a shop selling curtains cut to a
        customer-specified width.
  \item \textbf{(Short)} Why is checking \texttt{\$\_FILES['file']['type']}
        not a security control, and what is the server-side equivalent?
  \item \textbf{(Short)} An administrator raises the upload size limit in the
        admin panel and large files are still refused. Explain, and say where
        else to look.
  \item \textbf{(Applied)} A clothing shop has 300 styles, each in 5 sizes and
        4 colours. Compute the row counts for options and for variants across
        \texttt{oc\_product} and \texttt{oc\_product\_description}. Then make a
        recommendation, naming the single requirement that decides it.
  \item \textbf{(Applied)} Design the upload handling you would write for a
        shop that must accept customer artwork. State every check, say which
        side it runs on, say what each one defends against, and say what you
        would do with the file after accepting it so that serving it later is
        safe.
\end{reviewq}
